\documentclass[10pt]{article}

\usepackage[utf8]{inputenc}

\usepackage[T1]{fontenc}

\usepackage{amsmath}

\usepackage{amsfonts}

\usepackage{amssymb}

\usepackage[version=4]{mhchem}

\usepackage{stmaryrd}

\usepackage{graphicx}

\usepackage[export]{adjustbox}

\graphicspath{ {./images/} }



\begin{document}

ровании такая программа наз. и н т е $\mathrm{p}$ пा $\mathrm{p}$ е т а т ор о м входного языка.) Тогда справедливы следующие соотнонения:



\[

\begin{gathered}

\forall d N S(N L I, i p, d)=o p(d), \\

\forall i p N S(N S, N L I, i p)=t(i p), \\

\forall N L I N S(N S, N S, N L I)=t t(N L I),

\end{gathered}

\]



то есть объектная программа - это проекция интерпретатора входного языка ва входную программу; транслятор - әто гроекция частичного вычислителя на ин'ерпретатор входного языка, а транслятор трансляторов - это проекция частичного вычислителя на самото себя.



[2] JIиm.: [1] E p ш о в Ю. Л., Теория нумераций, М., 1977; [2] К Л и I и С. K., Введение в метаматематику, пер. с англ., M., 1957; [3] У c II е н c Is и it B. А., Јекции о вычислимых Фупкциях, М., 1960; [4] Е р ш о в А. ІІ., в кн.: Всесоюзная конференция. Методы математической логики в проблемах исПаланга, $3-5$ сент. 1980 , ч. 2, Вильнюс, 1980 , с. $26-55$.



ТРАНСМИССИИ УСЈОВИЕ - условие на налсевдодифђеренциальный оператор на гладком многообразии с краєм, гарантирующее, что гладкие вілоть до края функции, продолженные нулем, пөреводятся этим оператором снова в функции, гладкие вплоть до границы. Продолжение нулем здесь делается на некоторую окрестность исходного многообразия, которое считается вложенным в более широкое многообразие без края, так что точки грая становятся внутренними точками.



Если символ рассматриваемого псевдодифференциального оператора имеөт в локальных координатах в окрестности границы асимптотич. разложение по положительно однородным функциям $a_{\alpha}(x, \xi)(\alpha-$ порядок однородности), то Т. у. можно записать в виде следуюшего условия на функции $a_{\alpha}$ :



\[

\begin{aligned}

& \partial_{x}^{\gamma} \partial_{\xi^{\prime}}^{\beta}\left[a_{\alpha}\left(x, \xi^{\prime}, \xi_{n}\right)-\right. \\

& \left.-e^{-i \pi \alpha} a_{\alpha}\left(x,-\xi^{\prime},-\xi_{n}\right)\right]_{x_{n}=0, \xi^{\prime}=0}=0,

\end{aligned}

\]



где $\gamma, \beta$ - любые мультииндексы, координаты $x$ считаются выбранными в окрестности граничной точки так, что $\left\{x_{n}=0\right\}$ - уравнение края, $x_{n} \geqslant 0$ на самом многообразии, $\xi=\left(\xi^{\prime}, \xi_{n}\right)-$ координаты, дуальные к координатам $x=\left(x^{\prime}, x_{n}\right)$.



Jит.: [1] Ә с к и н Г. И., Краевые вадачи для эллиптических псевдодифференциальных уравнений, М., 1973; [2] В оu t e t d e M o n vel L., "Acta math.", 1971, v. 126, № 1/2,

p. $11-51$.



ТРАНСПОНИРОВАННАЯ МАТРИЦА - матрица, получающаяся из данной (прямоугольной или квадратной) матрицы $A=\left\|a_{i k}\right\|, i=1, \ldots ., m, k=1, \ldots, n$, после замены строк одноименными столбцами, т. е. матрица $\left\|a_{i k}^{\prime}\right\|$, где $a_{i k}^{\prime}=a_{k i}, i=1, \ldots, n, k=1, \ldots, m$. Число строк Т. м. равно числу столбцов матрицы $A$, а число столбцов - числу строк матрицы $A$. Матрицу, транспонированную по отнопению к матрице $A$, принято обовначать $A^{\top}$ или $A^{\prime}$ O. А. Пваиюва.



ТРАНСПОНИРОВАННБЕ УРАВНЕНИЯ - 1) Т. У. для линейных алгебраич. систем - уравнения, имеюпцие транспонированные матрицы.



\begin{enumerate}

  \setcounter{enumi}{1}

  \item Т. У. для линейных интегральных уравнений уравнения с ядрами $K(x, y)$ и $K(y, x)$.

\end{enumerate}



Т. у. наз. также с о Іо з н ы м и, а с с о ц и и р о-



\includegraphics[max width=\textwidth, center]{2023_07_09_a0cc4316308a0acf2514g-1}

н е н и м и.



ТРАНСПОРТНАЯ ЗАДАЧА - один из наиболее важных частных случаев общеӥ задачи линейного программирования. Содержательно Т. з. формулируется следуіошим образом.



Пусть в пунктах $A_{1}, A_{2}, . ., A_{n}$ производится некрый однородный продукт, пршчем объем производства этого продукта в пункте $A_{i}$ составляет $a_{i}$ единиг, $i=1$, ..., . Произведенныйі в пунктах производства продукт должен быть доставлен в пуєкты потребления $B_{1}, B_{2}, \ldots, B_{n}$, причем объем потребления в пункте $B_{j}$ составляет $b_{j}$ единиц продукта. Предполагается, что транспортировка готовой продукции возможна из любого пункта производства в любой пункт потребления и травспортные издержки, приходящиеся на перевозку единицы продукта из пункта $A_{i}$ в пункт $B_{j}$, cocтавляют $c_{i j}$ денежных единиц. Задача состоит в организации такого плана перевозок, при к-ром суммарные транспортные издержки были бы мивимальными.



Формально задача ставится следующим образом. Пусть $x_{i j}$ - количество продукта, перевозимого из пункта $A_{i}$ в пункт $B_{j}$. Требуется, огределить совокупность из $m n$ величин $x_{i j}$, удовлетворнгцци уеловиям



\[

\left.\begin{array}{c}

\sum_{j=1}^{n} x_{i j}=a_{i}, i=1,2, \ldots, m, \\

\sum_{i=1}^{m} x_{i j}=b_{j}, j=1,2, \ldots, n, \\

x_{i j} \geqslant 0,

\end{array}\right\}

\]



и обраццащих в минимум линейную форму



\[

L(X)=\sum_{i=1}^{m} \sum_{j=1}^{n} c_{i j} x_{i j}

\]



Группа ограничений (1) связана с тем обстоятельством, что объем вывезенного из каждого пункта прроизводства продукта в точности равен объему произведенного в этом пункте продукта, а объем ввезенного в пункт потребления продукта в точности соответствует его потребности. При әтих ограничениях необходимым и достаточным условием для разрешимости Т. з. является выполнение у с ло в и я б а л а в с а



\[

\sum_{i=1}^{m} a_{i}=\sum_{j=1}^{n} b_{j}

\]



Специфическими для Т. з. являтотся следующие два обстоятельства: а) каждое из перемепных $x_{i j}$ входит в два уравнения системы (1); б) все коэффициенты при переменных $x_{i j}$ принимают лишь два значения 0 или 1. Условия а) и б) позволили разработать для репения Т. 3. алгоритмы, существенно более простые, чем симплексный метод, к-рый является одним из основных методов решения задач линейного программирования.



Наиболее извествыми из этих алгоритмов являются метод потенциалов и т. н. венгерский метод. М е т о д п о т е в ц и а л о в - это метод последовательного улучшения плана (перевозок) с использованием второї теоремы двойственности для проверки оптимальности [1]. В е в г е р с к и й м е т о д - это метод последовательного построения допустимого плана, к-рый автоматически оказывается оптимальным. В основе венгерского алгоритма лежит метод чередующихся цепей [2].



Известны следующие два обобщения классич. Т. з., являющиеся отражением встречаюшихся на практике ситуаций. О т к р ы т а я м о де ль Т. з.- это Т. 3. с нарушенным условием баланса (2), что означает либо превышение объема производства над объемом потребления, либо ваоборот. Такая задача сводится к классич. Т. з. путем введения фиктивного пункта производства (или потребления) с мощностью производства (или потребления), равноӥ разности объемов производства и потребления.



Много индексные т ранс по р тные зад а ч и при сохранении общей проблемы минимизации транспортных издержек учитывают неоднородность груза (продуктов производства) и неоднородность транспортных средств.



За рубежом Т. з. иногда наз. п р о б л е м о й $\mathrm{X}$ и qк $о$ к а.



Лит.: [1] Г о ль т т е й н Е. Г., ЮО д и Н Д. Б., Задачіп линейного программкрования гранснортного тйіа, М., 1969;





\end{document}